% =====================================================================
%  Kirigamizer: An Algorithm Proposal
%  Designing kirigami (cut + fold) patterns for any surface.
%
%  Author: Emre Dayangac
%
%  Companion to origamizer_algorithms.tex. Kirigamizer mirrors the
%  Origamizer pipeline but is allowed to CUT: wherever Origamizer hides
%  surplus curvature by tucking, Kirigamizer may instead cut. This file
%  states (i) how cut locations are determined and (ii) which parts of
%  Origamizer are reused, then gives pseudocode for every stage.
%
%  Building blocks: E. D. Demaine & T. Tachi, "Origamizer", SoCG 2017;
%  A. Ghassaei, E. D. Demaine, N. Gershenfeld, GPU origami simulation,
%  7OSME 2018. The pseudocode and the cut-placement criterion are the
%  proposal's own contribution.
% =====================================================================
\documentclass[11pt]{article}

\usepackage[margin=1in]{geometry}
\usepackage{amsmath,amssymb}
\usepackage{xcolor}
\usepackage{enumitem}
\usepackage{booktabs}
\usepackage[ruled,vlined,linesnumbered]{algorithm2e}
\usepackage[hidelinks]{hyperref}

% ---- algorithm2e keyword setup (matches origamizer_algorithms.tex) ---
\SetKwInput{KwData}{Input}
\SetKwInput{KwResult}{Output}
\SetKwComment{Comment}{$\triangleright$\ }{}
\SetKwFor{ForEach}{for each}{do}{end}
\SetKw{KwRet}{return}
\SetKwFunction{Kirigamizer}{Kirigamizer}
\SetKwFunction{PlanCuts}{PlanCuts}
\SetKwFunction{Unfold}{SeamedUnfold}
\SetKwFunction{Route}{RouteSeams}
\SetKwFunction{Emit}{EmitPattern}
\SetKwFunction{Verify}{SimulateVerifyOptimize}
\SetKwFunction{Tuck}{OrigamizerTuck}
\SetKwFunction{Import}{ImportAndCondition}
\DontPrintSemicolon

% ---- shorthand ------------------------------------------------------
\newcommand{\eps}{\varepsilon}
\newcommand{\defect}{\delta}
\newcommand{\cut}{\mathcal{C}}
\newcommand{\crease}{\mathcal{K}}
\newcommand{\tuckset}{\Theta}
\newcommand{\Off}{\mathcal{O}_{\eps}}

\title{\textbf{Kirigamizer: An Algorithm Proposal}\\[2pt]
\large Designing Kirigami (Cut\,+\,Fold) Patterns for Any Surface}
\author{Emre Dayangac}
\date{}

\begin{document}
\maketitle

\noindent
\textbf{Proposal.} A paper fold is an isometry, so it preserves intrinsic
curvature: every interior point of a flat sheet keeps exactly $2\pi$ of
material around it, and folding can only \emph{redistribute} that angle,
never change the total at a point. This single fact bounds what origami can
do. At a convex vertex the target subtends \emph{less} than $2\pi$, leaving
a surplus that a fold can hide, precisely Origamizer's \emph{tuck}. At a
saddle vertex the target subtends \emph{more} than $2\pi$, which no interior
point of a flat sheet can supply, so the fold is impossible and
Origamizer can only \emph{approximate} it (split the vertex, $\eps$-extra).
Negative curvature is therefore not an incidental weakness but the
\emph{single, structural} place where origami's ``isometry on a disk''
assumption fails.

A \emph{cut} lifts exactly that assumption. Where the sheet is cut, a point
lies on the \emph{boundary}, not the interior, so it is no longer pinned to
$2\pi$: its two lips may overlap (a cone, $<2\pi$, for free) or splay open
with a gusset (a saddle, $>2\pi$). Cuts thus realize \emph{both} signs of
curvature, exactly and locally, supplying the one operation origami
lacks, and only there. I therefore propose \emph{Kirigamizer}: keep
Origamizer's pipeline unchanged, it is already strong everywhere except
this obstruction, and introduce cuts precisely where curvature is
negative or where tucking would be wasteful. The proposal contributes
(i) a criterion that \emph{determines where cuts must go}, derived from the
intrinsic curvature of $Q$, and (ii) a construction that \emph{reuses
Origamizer's modules} for everything cutting does not handle. Origamizer is
the special case $\cut=\varnothing$.
Symbols: $Q$ is the target (any genus $g$, any curvature, clean/tuck
sides); $\eps>0$ the tolerance; output is a flat sheet $P$ with cut set
$\cut$ and crease set $\crease$, and a folding $f\colon P\to\mathbb{R}^3$.

% =====================================================================
\section{Math}
% =====================================================================
Angle defect (discrete Gaussian curvature) and the curvature budget
(Gauss-Bonnet):
\begin{equation}
\defect(v)=2\pi-\sum_i\alpha_i(v),
\qquad
\sum_{v}\defect(v)=2\pi\,\chi(Q)=2\pi(2-2g).
\label{eq:gb}
\end{equation}
Flat paper carries exactly $2\pi$ at every interior point, so the local
action at a vertex is forced by the \emph{sign} of $\defect(v)$:
\begin{equation}
\text{at } v:\quad
\begin{cases}
\defect(v)>0 & \text{remove } \defect(v) \;\Rightarrow\; \textbf{dart (cut+overlap)} \text{ \emph{or} tuck},\\[2pt]
\defect(v)=0 & \text{nothing},\\[2pt]
\defect(v)<0 & \text{add } |\defect(v)| \;\Rightarrow\; \textbf{open slit \emph{or} gusset}.
\end{cases}
\label{eq:rule}
\end{equation}
Correctness target (Origamizer's containment, cuts allowed):
\begin{equation}
Q\subseteq f(P)\subseteq Q\cup\Off(Q)
\;\Longrightarrow\;
d_H\!\big(Q,f(P)\big)\le\eps,
\label{eq:target}
\end{equation}
with $\Off(Q)$ the tuck-side $\eps$-offset. A cut is \emph{sealed} when its
lips $c_1,c_2$ fold together, $d_F(f(c_1),f(c_2))\le\eps$; else it is left
\emph{open} (a deliberate hole).

% =====================================================================
\section{Why kirigami, not origami?}\label{sec:why}
% =====================================================================
The case for cutting is not aesthetic; it is forced by geometry, efficiency,
and topology.

\paragraph{(a) Origami can only \emph{subtract} angle; kirigami can also
\emph{add} it.} Rule \eqref{eq:rule} is asymmetric for a single uncut
sheet. At a convex vertex ($\defect>0$) the paper has surplus angle, which a
fold can hide (tuck). At a saddle vertex ($\defect<0$) the surface needs
\emph{more} than the $2\pi$ a flat interior point can supply, no fold can
create it. So origami covers only \emph{half} the curvature range; the
deficit half is reachable only by a cut (open slit or gusset). This is the
root reason, and it is exactly the operation Origamizer lacks.

\paragraph{(b) Negative curvature is the common case, not a corner case.}
Any target with a neck, a saddle, a handle, or genus $g>0$ has
$\defect<0$ vertices, and by \eqref{eq:gb} a genus-$g$ surface carries
total curvature $2\pi(2-2g)$, so for $g\ge 2$ the surface is \emph{net
negatively curved}. For all of these, an origami-only Origamizer must fall
back on its $\eps$-extra construction, which \emph{never folds the surface
exactly}: it splits each saddle vertex and approximates it within $\eps$.
Kirigami realizes those vertices \emph{exactly} with a cut.

\paragraph{(c) Cutting is dramatically more material-efficient.} On the
Stanford bunny ($374$ triangles) Origamizer puts only $\approx\!22.3\%$ of
the paper onto the visible surface, nearly $78\%$ is hidden tuck, a
$\sim\!2{:}1$ scale penalty per dimension. The tuck behind a convex vertex
grows with its surplus $\defect(v)$ and the size of the surrounding facets.
A \emph{dart} discards the same surplus through a thin slit whose lips
merely overlap by angle $\defect(v)$, no deep buried pleat. The saved
area scales with the total positive curvature $\sum_v\max(\defect(v),0)$,
which is large for any spiky or detailed model.

\paragraph{(d) Fewer layers $\Rightarrow$ actually foldable.} Tucks stack
many paper layers at each vertex; the bunny's pattern took a $\sim\!10$-hour
hand fold and is, in the authors' words, likely impossible to design by
hand. Cuts collapse those layer stacks, so the result is thinner and
foldable from stiff stock (cardstock, sheet metal, composites) where a
thick origami tuck simply cannot close.

\paragraph{(e) Topology and function come free with cuts.} A single uncut
sheet is a disk: holes, handles, and non-orientable targets must first be
\emph{cut} into a disk before Origamizer even starts. Kirigami takes those
cuts as part of the design. Moreover, whole classes of useful structures,
auxetic / negative-Poisson metamaterials, stretchable electronics,
deployable kirigami honeycombs, are \emph{defined} by their cuts and have
no origami-only realization at all.

\paragraph{Summary.} Origami is the special case $\cut=\varnothing$, optimal
only when the target is (nearly) developable and convex. The moment the
target has saddles, handles, fine detail, or a functional cut pattern,
i.e.\ most real targets, kirigami is not a stylistic choice but the
\emph{exact, efficient, manufacturable} one. The rest of this proposal makes
that switch precise: \S\ref{sec:where} decides where to cut, and
\S\ref{sec:reuse} keeps everything origami still does best.

% =====================================================================
\section{Input: conditioning and file format}\label{sec:input}
% =====================================================================
The whole method rests on the angle defect $\defect(v)=2\pi-\sum_i\alpha_i(v)$,
which is only well defined on a \emph{manifold, flat-faceted, coarse} mesh
with shared vertices. So the raw 3D model must first be conditioned, and the
choice of file format matters.

\paragraph{STL vs.\ STEP.} \emph{Prefer STEP (or a topological mesh: OBJ /
PLY / OFF); avoid raw STL.}
\begin{itemize}[leftmargin=1.4em,itemsep=2pt]
  \item \textbf{Connectivity.} $\defect(v)$ needs vertices \emph{shared}
  between faces. STL is triangle soup: duplicated vertices, no adjacency, so
  $\defect$ is undefined until the mesh is welded and its topology rebuilt.
  STEP is a B-rep: exact faces, edges, and adjacency are already present.
  \item \textbf{Feature edges.} Cut routing \eqref{eq:opt} wants to follow
  real edges (to stay seamless). STEP carries them; STL does not.
  \item \textbf{Units / scale.} The tolerance $\eps$ and any fabrication need
  true units. STEP stores units; STL is unitless and often inconsistent.
  \item \textbf{Curvature placement.} STEP surfaces are smooth (NURBS), so
  curvature is continuous; you must \emph{facet} them, which is exactly the
  step that decides where curvature concentrates (the cone vertices). On a
  dense STL the opposite problem bites: almost every vertex has
  $\defect\neq0$, so \eqref{eq:necessity} would cut \emph{everywhere} unless
  the mesh is first decimated to a coarse flat-faceted target.
\end{itemize}
\textbf{Decision: STEP is the input format.} It carries exact topology,
feature edges, and units, and lets us choose the faceting (hence the cone
vertices) ourselves, which is precisely what the cut planner needs. The
fixed pipeline is
\[
\textbf{STEP} \;\rightarrow\; \text{condition to coarse manifold mesh } Q
\;\rightarrow\; \text{half-edge internally}
\;\rightarrow\; \textbf{FKLD} \;+\; \textbf{SVG/DXF},
\]
where FKLD (a FOLD superset: folds, cuts, molecules, and material together,
\S\ref{sec:fold}) runs directly in the bar-and-hinge simulator and SVG/DXF
(cut vs.\ crease on separate layers) drives a cutter/plotter. A mesh (OBJ/PLY/OFF) is accepted as a fallback, and
raw STL only after weld $+$ manifold repair $+$ decimation.

\begin{algorithm}[H]
\caption[\textsc{ImportAndCondition}]{\textsc{ImportAndCondition}$(M)$: STEP model $\to$ usable target $Q$ (Stage 0)}
\KwData{STEP B-rep model $M$ (fallback: a mesh OBJ/PLY/OFF/STL).}
\KwResult{Coarse, manifold, flat-faceted target $Q$ with consistent units
          and clean/tuck sides.}
\BlankLine
\uIf{$M$ is STEP (B-rep)}{ read units and feature edges; tessellate each face into planar facets at the desired coarseness\;}
\Else(\Comment*[h]{fallback: $M$ is a mesh}){ weld coincident vertices; rebuild edge/face adjacency\;}
Repair to a 2-manifold (fill/zip cracks, fix orientation, remove degenerates)\;
Decimate / remesh to a \emph{coarse} flat-faceted model so curvature
  concentrates at a chosen vertex set (cone vertices)\;
  \Comment*[r]{otherwise $\defect\neq0$ almost everywhere $\Rightarrow$ \eqref{eq:necessity} cuts everywhere}
Record units; assign the clean side (visible) and tuck side (hidden)\;
\KwRet $Q$\;
\end{algorithm}

% =====================================================================
\section{Where the cuts go (the determination)}\label{sec:where}
% =====================================================================
Cut placement is \emph{derived}, not guessed, in three steps:
\emph{necessity}, \emph{connection}, \emph{routing}.

\paragraph{(1) Necessity: which vertices force a cut.}
A surface flattens isometrically to the plane iff it is a simply connected
\emph{developable disk}: no interior vertex with nonzero defect, and genus
$0$. Cutting moves curvature from the interior to a boundary, where it is
realized as an overlap (dart) or a gap (slit). Hence the hard constraint
\begin{equation}
v \in \cut \quad\Longleftrightarrow\quad \defect(v)\neq 0,
\label{eq:necessity}
\end{equation}
i.e.\ \emph{every defect vertex must become a boundary vertex of some
patch}. By \eqref{eq:gb} the \emph{total} signed curvature is fixed by
topology, so the amount of cutting is topologically forced, the proposal
only chooses \emph{how} to realize it. In addition, for genus $g>0$ a set
of $2g$ non-contractible loops $H$ must be cut to reduce the genus to $0$.

\paragraph{(2) Connection: making patches developable.}
Isolated punctures do not flatten; the cut set must be a \emph{graph}
$\cut$ such that every component of $Q\setminus\cut$ is a disk whose
interior is flat. Equivalently, $\cut$ must connect every defect vertex
(\ref{eq:necessity}) to a patch boundary, and include the handle loops
$H$. This is a Steiner-forest / cut-tree condition on the mesh:
\begin{equation}
\text{required terminals } \;=\; \{\,v : \defect(v)\neq 0\,\}\ \cup\ \partial Q,
\qquad \cut \supseteq H.
\end{equation}

\paragraph{(3) Routing: the exact location, by optimization.}
Among all $\cut$ satisfying (1)-(2), pick one minimizing a cost that
trades cut length against \emph{visibility} (cuts on the clean side are
penalized; cuts on edges of $Q$ or on the tuck side are cheap):
\begin{equation}
\min_{\cut}\ \sum_{e\in\cut}\Big(\ \mathrm{len}(e)\ +\ \lambda\,\mathrm{vis}(e)\ \Big)
\quad\text{s.t. \eqref{eq:necessity}, patches developable, } \cut\supseteq H,
\label{eq:opt}
\end{equation}
with $\lambda$ the seam-visibility weight. Practically: route $\cut$ along
\emph{mesh edges of $Q$} (so the clean side stays seamless), connect defect
vertices by \emph{shortest paths / geodesics} (minimizing length), and push
unavoidable interior cuts onto the \emph{tuck side}. Finally, the action at
each cut vertex follows the sign rule \eqref{eq:rule}; for $\defect(v)>0$
there is a free choice \emph{dart vs.\ tuck} (\S\ref{sec:reuse}), resolved
inside the optimization \eqref{eq:opt} by comparing a dart's seam cost
against an Origamizer tuck.

\begin{algorithm}[H]
\caption[\textsc{PlanCuts}]{\textsc{PlanCuts}$(Q,\lambda)$: decide where the cuts go (\S\ref{sec:where})}
\KwData{Surface $Q$; seam-visibility weight $\lambda$.}
\KwResult{Cut set $\cut$ (each cut tagged dart/slit/gusset) and tuck set
          $\tuckset$ of vertices to be handled by Origamizer.}
\BlankLine
\ForEach{vertex $v$ of $Q$}{ $\defect(v)\leftarrow 2\pi-\sum_i\alpha_i(v)$ }
$D\leftarrow\{v:\defect(v)\neq 0\}$\;
  \Comment*[r]{necessity \eqref{eq:necessity}: only these vertices force a cut}
$H\leftarrow$ $2g$ non-contractible loops of $Q$\;
  \Comment*[r]{cut handles so each patch has genus $0$}
$\cut\leftarrow$ min-cost graph on mesh edges linking $D\cup\partial Q$, plus $H$,
            minimizing \eqref{eq:opt}\;
  \Comment*[r]{prefer edges of $Q$; shortest paths; hide on the tuck side}
$\tuckset\leftarrow\varnothing$\;
\ForEach{$v\in D$}{
  \uIf{$\defect(v)<0$}{ tag $v$'s cut \emph{slit} (open) or \emph{gusset} of angle $|\defect(v)|$ }
  \uElseIf{$\defect(v)>0$ \textbf{and} darting is cheaper than tucking}{ tag $v$'s cut \emph{dart} of angle $\defect(v)$ }
  \Else(\Comment*[h]{$\defect(v)>0$, seamlessness preferred}){ move $v$ from $\cut$ to $\tuckset$ }
}
\KwRet $(\cut,\tuckset)$\;
\end{algorithm}

% =====================================================================
\section{What it reuses from Origamizer}\label{sec:reuse}
% =====================================================================
Kirigamizer is not a rewrite: every step that cutting does not handle is an
Origamizer module, invoked unchanged. Cutting is a \emph{front-end} that
removes exactly the curvature Origamizer struggles with; Origamizer then
does the rest.

\begin{center}
\renewcommand{\arraystretch}{1.25}
\begin{tabular}{@{}p{0.30\linewidth}p{0.62\linewidth}@{}}
\toprule
\textbf{Origamizer part} & \textbf{How Kirigamizer uses it} \\
\midrule
Angle-defect / curvature analysis (\S3) & drives cut placement \eqref{eq:necessity}-\eqref{eq:rule} instead of the tuck-proxy inset \\
Vertex-unfolding (\S4) & flattens each developable patch in \textsc{SeamedUnfold} \\
Tutte embedding (\S5) & packs the unfolded patches into one sheet $P$ \\
Streams: rivers \& brooks (\S5) & become \emph{seams/tabs} that rejoin lips, the same channels, run in reverse to \emph{deliver} material to $\defect<0$ vertices \\
Tuck proxy / tucking molecules (\S3,\S7) & invoked \emph{locally} for the $\defect>0$ vertices in $\tuckset$ we choose to tuck (keeps the clean side seamless there) \\
Site placement + Voronoi folding (\S6-\S7) & generates the actual crease pattern for each tucked region \\
Guarantees: watertight / seamless / $\eps$-extra & generalized to \emph{seam}-watertight / \emph{seam}-seamless / $\eps$-extra \eqref{eq:target} \\
\bottomrule
\end{tabular}
\end{center}

\noindent The division of labor: \textbf{cuts} take the cases Origamizer
cannot do exactly ($\defect<0$, genus, nonorientable); \textbf{Origamizer's
tucking} takes the cases it already does well ($\defect>0$ where a seamless
single sheet is wanted). The \emph{dart-vs-tuck} switch in \textsc{PlanCuts}
is exactly the dial between the two regimes.

% =====================================================================
\section{Top-level driver}
% =====================================================================
\begin{algorithm}[H]
\caption[\textsc{Kirigamizer}]{\textsc{Kirigamizer}$(M,\eps,\lambda)$: design a cut$+$fold pattern for model $M$}
\KwData{STEP model $M$ (B-rep); tolerance $\eps>0$; seam-visibility weight
        $\lambda$.}
\KwResult{Sheet $P$ with cuts $\cut$, creases $\crease$, folding $f$ with
          $f(P)\subseteq Q\cup\Off(Q)$.}
\BlankLine
$Q\leftarrow$ \Import{$M$}\;
  \Comment*[r]{Stage 0 (\S\ref{sec:input}): weld/repair/coarsen to a flat-faceted manifold}
$(\cut,\tuckset)\leftarrow$ \PlanCuts{$Q,\lambda$}\;
  \Comment*[r]{\S\ref{sec:where}: where the cuts go; what stays tucked}
$(\{U_k\},\crease)\leftarrow$ \Unfold{$Q,\cut$}\;
  \Comment*[r]{cut into developable patches, flatten isometrically}
$(P,\cut,\crease)\leftarrow$ \Route{$\{U_k\},\cut,\crease$}\;
  \Comment*[r]{pack to one sheet; route seams/tabs}
\Emit{$P,\crease,\cut,\tuckset$}\;
  \Comment*[r]{write crease$+$cut pattern; tuck $\tuckset$ via Origamizer}
$f\leftarrow$ \Verify{$P,\crease,\cut,Q,\eps$}\;
  \Comment*[r]{GPU-fold, check $d_H\le\eps$, refine}
\KwRet $(P,\cut,\crease,f)$\;
\end{algorithm}

% =====================================================================
\section{Stage algorithms}
% =====================================================================

\noindent\textbf{Stage 1: Plan cuts.} Given in \textsc{PlanCuts}
(\S\ref{sec:where}); replaces Origamizer's tuck proxy.

\BlankLine
\noindent\textbf{Stage 2: Seamed unfolding} (reuses vertex-unfolding).
With the cuts in place, $Q$ separates into developable patches that flatten
\emph{isometrically}, no $\eps$-approximation is needed for saddle
vertices, since their excess has been absorbed by a cut.

\begin{algorithm}[H]
\caption[\textsc{SeamedUnfold}]{\textsc{SeamedUnfold}$(Q,\cut)$: cut into patches and flatten}
\KwData{Surface $Q$; cut set $\cut$.}
\KwResult{Planar patches $\{U_k\}$; crease set $\crease$.}
\BlankLine
$\{Q_k\}\leftarrow$ connected components of $Q$ after cutting along $\cut$\;
$\crease\leftarrow\varnothing$\;
\ForEach{patch $Q_k$}{
  $U_k\leftarrow$ vertex-unfold $Q_k$ isometrically into the plane\;
    \Comment*[r]{developable $\Rightarrow$ isometric immersion exists}
  \While{$U_k$ self-overlaps}{
    Add a relief cut across $Q_k$ and re-split; re-unfold\;
    \Comment*[r]{an immersion need not embed; relief cuts restore an embedding}
  }
  Add the interior fold edges of $U_k$ (with M/V assignment) to $\crease$\;
}
\KwRet $(\{U_k\},\crease)$\;
\end{algorithm}

\noindent\textbf{Stage 3: Seam \& tab routing} (reuses streams). Pack the
patches into one sheet and connect the lips that must rejoin: a zero-width
\emph{seam} (brook) for sealed cuts, or a thickened \emph{tab/gusset strip}
(river) that carries extra material to $\defect<0$ vertices; open cuts stay
as boundaries of $P$.

\begin{algorithm}[H]
\caption[\textsc{RouteSeams}]{\textsc{RouteSeams}$(\{U_k\},\cut,\crease)$: pack and connect}
\KwData{Patches $\{U_k\}$; cut set $\cut$; crease set $\crease$.}
\KwResult{Sheet $P$; updated $\cut,\crease$.}
\BlankLine
$P\leftarrow$ pack $\{U_k\}$ into one sheet (Tutte-embed patch-adjacency graph)\;
\ForEach{cut $c\in\cut$ marked \emph{sealed}}{
  \uIf{$c$ adds angle (a $\defect<0$ vertex)}{ route a \emph{tab/gusset strip} (river) between the lips; add folds to $\crease$ }
  \Else{ route a zero-width \emph{seam} (brook) matching the lips }
}
Leave every \emph{open} cut as a free boundary of $P$\;
\KwRet $(P,\cut,\crease)$\;
\end{algorithm}

\noindent\textbf{Stage 4: Emit pattern} (reuses site placement \&
Voronoi folding). Write mountain/valley creases plus cut edges in the
\textbf{FKLD} format (\S\ref{sec:fold}); for each tucked vertex, generate its
Origamizer tuck locally and merge.

\begin{algorithm}[H]
\caption[\textsc{EmitPattern}]{\textsc{EmitPattern}$(P,\crease,\cut,\tuckset)$: write the crease$+$cut pattern}
\KwData{Sheet $P$; crease set $\crease$; cut set $\cut$; tuck set $\tuckset$.}
\KwResult{An FKLD file over $P$ (\S\ref{sec:fold}).}
\BlankLine
\ForEach{edge $k\in\crease$}{ mark $k$ mountain/valley/facet with target fold angle }
\ForEach{edge $c\in\cut$}{ mark $c$ as a cut (boundary) edge }
\ForEach{vertex $v\in\tuckset$}{
  $\crease\leftarrow\crease\ \cup\ $\Tuck{$v$}\;
    \Comment*[r]{Origamizer \S6-\S7: local Voronoi tuck (geodesic-vertex sites)}
}
Export $(P,\crease,\cut)$\;
\end{algorithm}

\noindent\textbf{Stage 5: Simulate, verify, optimize} (replaces the
constructive proof). The GPU bar-and-hinge solver treats cuts as
first-class edges (\S\ref{sec:sim}: full cuts split the mesh, partial cuts
soften the hinge by $(1-\mathrm{cutRatio})$), and a driven-boundary forward
process steers it to the goal mesh $M_0$; iterate until the folded state lies
within $\eps$ of $Q$.

\begin{algorithm}[H]
\caption[\textsc{SimulateVerifyOptimize}]{\textsc{SimulateVerifyOptimize}$(P,\crease,\cut,Q,\eps)$}
\KwData{Pattern $(P,\crease,\cut)$; target $Q$; tolerance $\eps$.}
\KwResult{Folding $f\colon P\to\mathbb{R}^3$ with $d_H(Q,f(P))\le\eps$.}
\BlankLine
\Repeat{$r\le\eps$}{
  $f\leftarrow$ \textsc{FoldSolver.solve}$(P,M_0,1,\mathrm{iters})$\;
    \Comment*[r]{AKDE explicit solver; driven boundary to $M_0$ (\S\ref{sec:sim})}
  $r\leftarrow d_H\!\big(Q,f(P)\big)$\;
  \If{$r>\eps$ \textbf{or} strain/self-intersection too high}{
    Adjust: cut routing $\lambda$, dart-vs-tuck in $\tuckset$, fold angles, scale\;
  }
}
\KwRet $f$\;
\end{algorithm}

\noindent\emph{Correctness target.}
\textbf{Seam-seamless:} each facet of $Q$ is one uncut face on the clean
side (cuts on edges / tuck side). \textbf{Seam-watertight:} every cut is
lip-matched ($d_F\le\eps$) or tuck-hidden, so the clean side is gap-free.
\textbf{$\eps$-extra:} $f(P)\subseteq Q\cup\Off(Q)$, hence $d_H\le\eps$
\eqref{eq:target}.

% =====================================================================
\section{Simulation: AKDE bar-and-hinge with driven boundary}\label{sec:sim}
% =====================================================================
Stage 5 uses the simulator from \textbf{AKDE} (\texttt{kirigami/sim}): the
explicit \emph{bar-and-hinge} solver of Ghassaei, Demaine, and Gershenfeld,
extended for kirigami. Its CPU force math (\texttt{forces.ts}) is the
unit-testable twin of the GPU GLSL passes; the mesh is stored struct-of-arrays
(one texel per node/edge/face). Masses are unit ($m\equiv1$). AKDE defaults:
$EA=20$, $k_{\text{fold}}=k_{\text{facet}}=0.7$, $k_{\text{face}}=0.2$,
damping $\zeta\approx0.45$ (AKDE runs overdamped, $\zeta\!\approx\!1$, to settle
the flat$\to$cone buckling), design angles $\theta_M\!\approx\!1.2$,
$\theta_V\!\approx\!2.9$ (valley near $\pi$ so each molecule tucks nearly flat,
collapsing the flat-net perimeter radius $s$ onto the base radius $R$).

\paragraph{Signed dihedral fold angle (branch-safe).} With $\hat e$ the unit
crease-edge vector and $n_1,n_2$ the adjacent face normals,
\begin{equation}
\theta=\operatorname{atan2}\!\big((n_1\times\hat e)\cdot n_2,\ n_1\cdot n_2\big),
\qquad \gamma=\pi-\theta,
\label{eq:theta}
\end{equation}
unwrapped against the previous step so it tracks continuously through $\pm\pi$
($\gamma$ is the geometric dihedral). This signed \texttt{atan2} form (DETC2019
Eq.\ 5) stays valid through the full fold range, unlike a $\sin\gamma$ closed form.

\paragraph{Axial bar $+$ neighbour damping (per edge).} With $\hat d$ the unit
edge vector, length $l$, rest $l_0$, $k_a=EA/l_0$:
\begin{equation}
F_{\text{ax}}=k_a\,(l-l_0)\,\hat d,\qquad
F_{\text{damp}}=c\,(v_b-v_a),\quad c=2\zeta\sqrt{k_a\,m_{\min}},
\end{equation}
i.e.\ damping is \emph{relative} between the two endpoints, not a global drag.

\paragraph{Crease torque on four nodes.} For the two wing vertices $w\in\{1,2\}$
(opposite the crease, with face normal $\hat n_w$ and moment arm $h_w$ to the
crease line) and the two crease nodes split by along-crease fractions
$\mathrm{coef}_w$, scaled by the fold slider $\mathrm{fp}\in[0,1]$:
\begin{equation}
\Phi=k_c\,(\mathrm{fp}\cdot\theta_{\text{tgt}}-\theta),\qquad
F_{\text{wing }w}=\frac{\Phi}{h_w}\,\hat n_w,
\end{equation}
with momentum-conserving reactions on the crease nodes weighted by
$(1-\mathrm{coef}_w)$ and $\mathrm{coef}_w$.

\paragraph{Kirigami coupling (DETC2019-97557, Eq.\ 6: $Kd=Fx$).} A cut softens
the hinge in proportion to lost ligament:
\begin{equation}
k_c=(1-\mathrm{cutRatio})\,l_0\,k,\qquad k\in\{k_{\text{fold}},k_{\text{facet}}\},
\label{eq:cut}
\end{equation}
so $\mathrm{cutRatio}=0$ is exactly Gershenfeld and $\mathrm{cutRatio}=1$ is a
free flap. Fully cut (major/minor) edges are simply \emph{absent as creases}:
their lips are independent nodes carrying only a boundary bar, which is what
lets holes and darts open. (This is AKDE's actual mechanism, replacing the
earlier ``seam spring'' sketch.)

\paragraph{Face interior-angle spring.} For each triangle angle $\alpha$ with
flat value $\alpha_0$, a torque from $k_{\text{face}}(\alpha_0-\alpha)$ applied
via the gradient $\operatorname{cross}(n,e)/|e|^2$ on the three nodes, resisting
in-plane shear.

\paragraph{Explicit integration.} Forward Euler at unit mass, with the stable
step bounded by the stiffest (axial) mode:
\begin{equation}
v \leftarrow v+\tfrac{F}{m}\,\Delta t,\qquad p \leftarrow p+v\,\Delta t,\qquad
\Delta t=\frac{0.9}{2\pi\,\omega_{\max}},\ \ \omega_{\max}=\max\sqrt{k_a/m_{\min}}.
\label{eq:euler}
\end{equation}
State lives in GPU textures; one thread per element, so per-step cost is roughly
constant until node count exceeds core count, then linear.

\paragraph{Driven-boundary forward process (the design coupling, DETC \S3.2).}
This is how the simulator is \emph{steered} to the designed shape, not merely
checked against it. Boundary nodes are driven kinematically toward the goal mesh
$M_0$ (the folded target $Q$) and pinned, so the force passes leave them in place
while the interior relaxes:
\begin{equation}
p_i=\mathrm{rest}_i+(\mathrm{goal}_i-\mathrm{rest}_i)\cdot\mathrm{fp},
\qquad i\in\text{driven}.
\label{eq:drive}
\end{equation}
The slider $\mathrm{fp}$ is eased \emph{quasi-statically},
$\mathrm{fp}\leftarrow\mathrm{fp}+(\mathrm{target}-\mathrm{fp})\cdot0.004$ per
step, because a fast ramp injects energy the explicit integrator cannot
dissipate.

\begin{algorithm}[H]
\caption[\textsc{FoldSolver}]{\textsc{FoldSolver.solve}$(P,M_0,\mathrm{target},\mathrm{iters})$: AKDE explicit kirigami solver (\S\ref{sec:sim})}
\KwData{Bar-hinge model of $P$ (beams, creases with $\theta_{\text{tgt}}$ and $\mathrm{cutRatio}$, faces); goal mesh $M_0$; target fold fraction; step count.}
\KwResult{Folded positions $f$ (boundary driven to $M_0$, interior relaxed).}
\BlankLine
Build model: split mesh at major/minor cuts; set $k_c$ by \eqref{eq:cut}; mark driven boundary nodes\;
$\Delta t\leftarrow 0.9/(2\pi\,\omega_{\max})$; $\ \mathrm{fp}\leftarrow 0$; $\ v\leftarrow 0$\;
\For{$\mathrm{iters}$ steps}{
  $\mathrm{fp}\leftarrow \mathrm{fp}+(\mathrm{target}-\mathrm{fp})\cdot 0.004$ \Comment*[r]{quasi-static ease}
  drive boundary by \eqref{eq:drive} for $i\in\text{driven}$\;
  pass 1: face normals $n_t$\;
  pass 2: fold angles $\theta_e$ \eqref{eq:theta} (unwrapped)\;
  pass 3: $F\leftarrow F_{\text{ax}}+F_{\text{damp}}+F_{\text{crease}}+F_{\text{face}}$\;
  pass 4: integrate \eqref{eq:euler}; skip fixed/driven nodes\;
}
\KwRet $f\leftarrow p$\;
\end{algorithm}

% =====================================================================
\section{FKLD output format}\label{sec:fold}
% =====================================================================
\textsc{EmitPattern} writes \textbf{FKLD} (Flexible Kirigami list
Datastructure), AKDE's format: a strict \emph{superset} of FOLD 1.2 (Demaine,
Ku, Lang). Every FKLD file is a valid FOLD file, so any FOLD tool renders the
crease pattern and ignores the kirigami extensions, which live under the
reserved \texttt{fkld:} namespace. The standard FOLD arrays carry the geometry:
\texttt{vertices\_coords}, \texttt{edges\_vertices}, \texttt{edges\_assignment}
(\texttt{M}/\texttt{V}/\texttt{B}/\texttt{U}/\texttt{F} plus \texttt{C} = cut),
\texttt{edges\_foldAngle}, \texttt{faces\_vertices}, \texttt{faceOrders} (tuck
stacking), and \texttt{file\_frames} (flat / target / folded states).

The kirigami extensions (FKLD v0.1 registry):
\begin{itemize}[leftmargin=1.4em,itemsep=1pt]
  \item \texttt{fkld:edges\_cutType}: subtype of each \texttt{C} edge, one of
  \texttt{major}, \texttt{minor}, \texttt{seam}, \texttt{dart},
  \texttt{auxetic}, \texttt{vent}, \texttt{tab}.
  \item \texttt{fkld:edges\_moleculeTheta} / \texttt{\_moleculeWidth} /
  \texttt{\_moleculeDepth}: tucking-molecule angle $\theta(i,j)$, chord width
  $w(i,j)$, and tuck depth behind the face (Tachi 2010).
  \item \texttt{fkld:edges\_dihedralTarget}: folded dihedral target (rad), i.e.\
  the $\theta_{\text{tgt}}$ consumed by the solver \eqref{eq:theta}.
  \item \texttt{fkld:vertices\_curvatureClass}
  (\texttt{positive}/\texttt{negative}/\texttt{flat}/\texttt{boundary}),
  \texttt{fkld:vertices\_angleDefect} $G_v=2\pi-\sum\theta$ \eqref{eq:gb}, and
  \texttt{fkld:vertices\_reliefStrategy}
  (\texttt{molecule}/\texttt{cut}/\texttt{hybrid}/\texttt{none}).
  \item \texttt{fkld:faces\_materialId} / \texttt{\_thickness} /
  \texttt{\_structuralRole} / \texttt{\_panelId} for fabrication.
\end{itemize}
The correspondences close the loop: \texttt{vertices\_reliefStrategy} \emph{is}
the dart-vs-tuck dial of \textsc{PlanCuts} (\S\ref{sec:where});
\texttt{vertices\_angleDefect} is \eqref{eq:gb}; the molecule keys carry the
Origamizer tuck (\S\ref{sec:reuse}); \texttt{edges\_dihedralTarget} feeds the
solver. For plain-FOLD tools, \texttt{filter.splitCuts} downgrades \texttt{C} to
\texttt{B} (lossy but legal; the \texttt{fkld:} annotations are preserved). A
schematic FKLD for one darted positive-curvature vertex:
\begin{verbatim}
{ "file_spec":1.2, "file_creator":"AKDE/Kirigamizer", "unit":"mm",
  "file_classes":["creasePattern"],
  "vertices_coords":[[0,0],[10,0],[10,10],[0,10],[5,5]],
  "edges_vertices":[[0,1],[1,2],[2,3],[3,0],[4,1],[4,3]],
  "edges_assignment":["B","B","B","B","V","C"],
  "edges_foldAngle":[0,0,0,0,150,0],
  "faces_vertices":[[0,1,4],[1,2,3,4],[3,0,4]],
  "fkld:edges_cutType":[null,null,null,null,null,"dart"],
  "fkld:vertices_curvatureClass":
       ["boundary","boundary","boundary","boundary","positive"],
  "fkld:vertices_angleDefect":[0,0,0,0,1.05],
  "fkld:vertices_reliefStrategy":
       ["none","none","none","none","hybrid"] }
\end{verbatim}

% =====================================================================
\section{Preconditions and limitations (honest scope)}\label{sec:limits}
% =====================================================================
What the pipeline does and does not guarantee:
\begin{itemize}[leftmargin=1.4em,itemsep=2pt]
  \item \textbf{Coarse flat-faceted input required.} The defect rule
  \eqref{eq:necessity} assumes curvature lives at vertices. The STEP model is
  faceted and conditioned to a coarse mesh in Stage 0 (\S\ref{sec:input});
  the faceting coarseness sets where curvature concentrates, and the output
  quality inherits that choice.
  \item \textbf{Open kirigami is fully general; sealed kirigami is not (yet).}
  Negative-curvature vertices are always realizable as \emph{open} slits
  (deliberate holes). The \emph{sealed} route (a gusset that delivers
  material while staying watertight) has no proof of a valid, non-crossing
  layer ordering in general; it is heuristic, checked in Stage 5.
  \item \textbf{Unfolding may need extra cuts.} A developable patch immerses
  isometrically but need not embed; the relief loop in \textsc{SeamedUnfold}
  adds cuts until it embeds, so the final cut count can exceed the curvature
  minimum.
  \item \textbf{No constructive flat-foldability guarantee.} Mountain/valley
  assignment and a collision-free layer order are produced heuristically and
  validated by the simulator, not proven (unlike Origamizer \S7).
  \item \textbf{Stage 5 is a search, not a solver.} The verify/refine loop
  has no convergence guarantee; it terminates when $d_H\le\eps$ or on a
  budget.
  \item \textbf{Zero-thickness model.} Material thickness, laser kerf, fold
  allowance/score lines, and joining tabs for sealed seams are out of scope
  here and must be added for real fabrication.
\end{itemize}
\noindent\textbf{Bottom line.} As written, the method reliably turns a
conditioned, coarse 3D model into an \emph{open}-kirigami pattern; the
watertight \emph{sealed} variant and guaranteed flat-foldability are the open
parts, deferred to the simulator and to future work.

\end{document}
